

\section{Diseño del control regulatorio (Nivel 2)}

El Nivel 2 corresponde al control regulatorio de las variables continuas del sistema. Su objetivo es garantizar que los 
accionamientos sigan referencias de posición, velocidad y aceleración con precisión y suavidad, respetando restricciones 
físicas (límites de recorrido, saturaciones de torque, corrientes máximas) y atenuando perturbaciones como el balanceo de 
la carga o variaciones de masa. Este nivel operará en tiempo discreto mediante muestreo (Sampled CTDS) y se implementará 
con lazos de realimentación en cascada (corriente/torque, velocidad y posición), además de compensaciones para mejorar 
la respuesta transitoria.

En este proyecto, el control regulatorio se encarga de los subsistemas de traslación del carro y de izaje,
coordinando sus dinámicas para lograr movimientos simultáneos y seguros. Las señales generadas en este nivel
son consignas y comandos de accionamiento que luego son habilitados o bloqueados por el control supervisor (Nivel 1) y,
ante eventos críticos, anulados por el nivel de protección (Nivel 0). De este modo, el Nivel 2 provee el desempeño
continuo requerido, mientras los niveles superiores garantizan la lógica de operación y la seguridad global del sistema.

\subsection{Control de Izaje de la Carga}

Para el control de la carga en el izaje, se diseñó una arquitectura de Control Directo de Velocidad
con Planificación Cinemática (Trayectoria $\to$ Velocidad $\to$ Torque). Esta estrategia evita los
sobresesfuerzos mecánicos, disminuye el chattering y optimiza los tiempos de ciclo.

\subsubsection{Modelo de Izaje para Control de Velocidad}

Partiendo de la ecuación dinámica del izaje y considerando el sistema linealizado alrededor de
un punto de operación, se asume que la dinámica del cable es rígida.

La ecuación de balance de torques en el eje del motor es:

\begin{equation}
    J_{\text{eq},h}\,\frac{d\omega_{hm}(t)}{dt} + b_{\text{eq},h}\,\omega_{hm}(t) = T_{hm}(t) - T_{\text{dist}}(t)
    \label{eq:din_izaje}
\end{equation}

donde $T_{\text{dist}}(t)$ es el torque perturbador que encapsula todas las dinámicas no lineales y externas
al accionamiento:

\begin{equation}
    T_{\text{dist}}(t) = \frac{F_{hw}(t)\,r_{hd}}{i_h} + T_{fb}(t, \text{BRK}_h) + \frac{T_{hEb}(t, \text{BRK}_{h,E})}
    {i_h}
    \label{eq:torque_dist}
\end{equation}

\noindent
\begin{tabular}[t]{@{}l@{\quad}l@{}}
\textit{Tensión de Cable} & $\frac{F_{hw}(t)\,r_{hd}}{i_h}$ \\
\textit{Fricción de Frenos} & $T_{fb}(t, \text{BRK}_h)$
\end{tabular}

Aplicando la Transformada de Laplace, la función de transferencia de la planta es:

\begin{equation}
    G_{\text{ved}}(s) = \frac{\Omega_{hm}(s)}{T_{hm}(s)} = \frac{1}{J_{\text{eq},h}\,s + b_{\text{eq},h}}
    \label{eq:FT_izaje}
\end{equation}

El polo mecánico del sistema es $P_{\text{mec}} = -\frac{b_{\text{eq},h}}{J_{\text{eq},h}}$. Dado que la fricción es muy baja en estos
sistemas, el polo se encuentra muy cerca del origen ($P_{\text{mec}} \approx 0$), comportándose casi como un
integrador puro.

\subsection{Identificación de Parámetros}

De manera análoga al carro, se realiza identificación experimental mediante respuesta al escalón de torque
en lazo abierto, extrayendo la constante de tiempo $\tau$ y la ganancia estática $K$ de la planta, que permiten
calcular:

\begin{align}
    b_{\text{eq},h} &= \frac{i_h\,T_0}{\omega_{\infty} - \omega_0} \\[4pt]
    J_{\text{eq},h} &= \tau\,b_{\text{eq},h}
\end{align}

\subsection{Diseño del Controlador PI para Izaje}

Se adopta una estrategia PI análoga a la del carro, con colocación de polos doble real en
$s = -\omega_{\text{cl},h}$ para garantizar respuesta sin sobreoscilación:

\begin{align}
    K_p = \frac{2\,\omega_{\text{cl},h}\,J_{\text{eq},h} - b_{\text{eq},h}}{i_h}
    \label{eq:Kp_izaje} \\[4pt]
    K_i = \frac{\omega_{\text{cl},h}^2\,J_{\text{eq},h}}{i_h}
    \label{eq:Ki_izaje}
\end{align}

La frecuencia de corte en lazo cerrado $\omega_{\text{cl},h}$ se elige como $\omega_{\text{cl},h} = \alpha\,|p_{\text{OL}}|$,
con $\alpha \in [3, 20]$ según los requisitos de velocidad de respuesta y robustez ante incertidumbres
de carga variable.

\begin{align}
\frac{B_{eq}*K_d*i_h}{J_{eq}} &= (2 \cdot \zeta + 1) \cdot \omega_{pos} \\
\frac{K_{p}*i_h}{J_{eq}} &= (2 \cdot \zeta + 1) \cdot \omega_{pos}^2 \\
\frac{K_{i}*i_h}{J_{eq}} &= \omega_{pos}^3
\end{align}

% \[
% \dot{\omega}=(B_{eq}\omega+T_{hm})/J_{eq}
% \]

% \[
% \dot{\theta}=\omega
% \]

\section{Identificación de Parámetros de la Planta e Implementación del Controlador PI}

\subsection{Modelo de la Planta de Velocidad de Izaje}

El accionamiento de izaje puede representarse, referido al eje rápido del motor,
mediante la ecuación diferencial de primer orden:

\begin{equation}
    J_{eq}\,\dot{\omega}_{hm}(t) = i_h\,T_{hm}(t) - B_{eq}\,\omega_{hm}(t)
\end{equation}

donde los parámetros equivalentes referidos al eje del motor son:

\begin{align}
    J_{eq} &= J_{hm+hb} + \frac{J_{hd+hEb}}{i_h^2}  \label{eq:Jeq}\\[4pt]
    B_{eq} &= b_{hm}     + \frac{b_{hd}}{i_h^2}       \label{eq:Beq}
\end{align}

Tomando la Transformada de Laplace de \eqref{eq:dinamica_motor} con condiciones
iniciales nulas, la función de transferencia de la planta resulta:

\begin{equation}
    G(s) = \frac{\omega_{hm}(s)}{T_{hm}(s)}
         = \frac{i_h}{J_{eq}\,s + B_{eq}}
    \label{eq:FT_planta}
\end{equation}

Expresada en forma estándar de primer orden:

\begin{equation}
    G(s) = \frac{K}{\tau\,s + 1},
    \qquad
    K = \frac{i_h}{B_{eq}},
    \qquad
    \tau = \frac{J_{eq}}{B_{eq}}
    \label{eq:FT_estandar}
\end{equation}

donde $K$ es la ganancia estática de la planta y $\tau$ su constante de tiempo.

% ----------------------------------------------------------------
\subsection{Identificación Experimental por Respuesta al Escalón}

Dado que los parámetros nominales calculados con \eqref{eq:Jeq}--\eqref{eq:Beq}
no contemplan efectos de carga variable ni la componente gravitacional, se adopta
el método de identificación por \textit{Process Reaction Curve}: se aplica un
escalón de torque conocido $T_0$ en lazo abierto y se registra la respuesta
$\omega_{hm}(t)$.

\subsubsection{Solución analítica ante escalón}

Para una entrada $T_{hm}(s) = T_0/s$, la salida en el dominio de Laplace es:

\begin{equation}
    \omega_{hm}(s)
    = G(s)\,\frac{T_0}{s}
    = \frac{K\,T_0}{s\,(\tau s + 1)}
\end{equation}

Antitransformando:

\begin{equation}
    \omega_{hm}(t)
    = \omega_0
      + \bigl(\omega_\infty - \omega_0\bigr)
        \!\left(1 - e^{-t/\tau}\right)
    \label{eq:respuesta_escalon}
\end{equation}

donde $\omega_0 = \omega_{hm}(0)$ es la condición inicial.

\subsubsection{Valor final por el Teorema del Valor Final}

Aplicando el Teorema del Valor Final a la salida:

\begin{equation}
    \omega_\infty
    = \lim_{s \to 0}\, s\,\omega_{hm}(s)
    = \lim_{s \to 0}\, s \cdot \frac{K}{\tau s + 1} \cdot \frac{T_0}{s}
    = K\,T_0
    = \frac{i_h\,T_0}{B_{eq}}
    \label{eq:valor_final}
\end{equation}

En régimen permanente $\dot{\omega} = 0$, de modo que \eqref{eq:valor_final}
equivale directamente al equilibrio estático de \eqref{eq:dinamica_motor}.

\subsubsection{Extracción de parámetros desde el ensayo}

A partir de la curva de respuesta se identifican experimentalmente:

\begin{itemize}
    \item $\omega_0$: valor de $\omega_{hm}$ en $t = 0$.
    \item $\omega_\infty$: valor en régimen permanente.
    \item $\tau$: tiempo en alcanzar
          $\omega_0 + 0{,}632\,(\omega_\infty - \omega_0)$,
          según \eqref{eq:respuesta_escalon}.
\end{itemize}

Los parámetros físicos equivalentes resultan:

\begin{align}
    B_{eq} &= \frac{i_h\,T_0}{\omega_\infty - \omega_0}
    \label{eq:Beq_id}\\[6pt]
    J_{eq} &= \tau\,B_{eq}
    \label{eq:Jeq_id}\\[6pt]
    p_{OL} &= -\frac{B_{eq}}{J_{eq}} = -\frac{1}{\tau}
    \label{eq:polo_OL}
\end{align}

% ----------------------------------------------------------------
\subsection{Diseño del Controlador PI por Colocación de Polos}

Se adopta un controlador Proporcional-Integral (PI):

\begin{equation}
    C(s) = K_p + \frac{K_i}{s}
         = \frac{K_p\,s + K_i}{s}
    \label{eq:PI}
\end{equation}

\subsubsection{Ecuación característica en lazo cerrado}

Con realimentación unitaria, los polos en lazo cerrado son las raíces de
$1 + C(s)\,G(s) = 0$:

\begin{equation}
    s\,(J_{eq}\,s + B_{eq})
    + i_h\,(K_p\,s + K_i) = 0
\end{equation}

Dividiendo por $J_{eq}$:

\begin{equation}
    s^2
    + \frac{B_{eq} + i_h\,K_p}{J_{eq}}\,s
    + \frac{i_h\,K_i}{J_{eq}}
    = 0
    \label{eq:ec_caract}
\end{equation}

\subsubsection{Polinomio deseado}

Se especifica un polo doble real en $s = -\omega_{cl}$
(factor de amortiguamiento $\zeta = 1$, sin sobreoscilación):

\begin{equation}
    (s + \omega_{cl})^2
    = s^2 + 2\,\omega_{cl}\,s + \omega_{cl}^2
    = 0
    \label{eq:pol_deseado}
\end{equation}

donde $\omega_{cl}$ es el único parámetro libre de diseño, elegido como:

\begin{equation}
    \omega_{cl} = \alpha\,|p_{OL}| = \frac{\alpha}{\tau},
    \qquad \alpha \in [3,\,20]
    \label{eq:wcl}
\end{equation}

El tiempo de estabilización al 2\% resulta aproximadamente:

\begin{equation}
    t_s \approx \frac{4}{\omega_{cl}}
\end{equation}

\subsubsection{Igualación de coeficientes}

Igualando \eqref{eq:ec_caract} con \eqref{eq:pol_deseado} término a término:

\begin{align}
    \frac{B_{eq} + i_h\,K_p}{J_{eq}} &= 2\,\omega_{cl}
    \label{eq:ec1}\\[4pt]
    \frac{i_h\,K_i}{J_{eq}} &= \omega_{cl}^2
    \label{eq:ec2}
\end{align}

Despejando las ganancias del controlador:

\begin{align}
    \boxed{K_p = \frac{2\,\omega_{cl}\,J_{eq} - B_{eq}}{i_h}}
    \label{eq:Kp}\\[8pt]
    \boxed{K_i = \frac{\omega_{cl}^2\,J_{eq}}{i_h}}
    \label{eq:Ki}
\end{align}

\subsubsection{Anti-windup por back-calculation}

Para evitar la saturación del integrador cuando la consigna de torque supera
$T_{hm,\text{MAX}}$ (límite físico del accionamiento), se implementa la corrección
por \textit{back-calculation}:

\begin{equation}
    \dot{x}_i(t) = e(t) + K_b\,\bigl[T^*_{sat}(t) - T^*(t)\bigr],
    \qquad
    K_b = \frac{1}{K_p}
\end{equation}

donde $T^*_{sat}$ es la salida del limitador y $T^*$ la salida no saturada del PI.
Cuando no hay saturación ambas coinciden y la corrección es nula.
Cuando el actuador satura, el término de corrección detiene la integración,
eliminando el fenómeno de \textit{windup}.


